% !PRINTED 206
\weberpage{206}

Der Hauptnenner der in (1) vorkommenden Brüche:

\[
\prod^{\nu} \vartheta_{g_1,g_2}\left(u + \frac{2\nu}{n}\right) = \pm \prod^{\nu} \vartheta_{g_1,g_2}\left(u + \frac{\nu}{n}\right)
\]

ist nach den letzten Formeln des § 33, von einem konstanten
Faktor abgesehen, gleich

\[
\vartheta_{g_1,g_2}(nu, n\omega),
\]

und wenn wir also die Summe (1) gleich

\weq{2}{\frac{\Phi(u)}{\vartheta_{g_1,g_2}(nu, n\omega)}}

setzen, so ist $\Phi(u)$ eine \gesperrt{ganze} transzendente Funktion von $u$.
Nun ist

\[
\frac{\vartheta_{11}\left(u + \frac{2\nu+1}{u}\right)}{\vartheta_{g_1,g_2}\left(u + \frac{2\nu+1}{u}\right)} = (-1)^{g_1+1} \frac{\vartheta_{11}\left(u + \frac{2\nu+n+1}{n}\right)}{\vartheta_{g_1,g_2}\left(u + \frac{2\nu+n+1}{n}\right)}
\]
% [sic] The source prints u in both left-hand inner denominators; n is intended.

und $\nu$ durchläuft gleichzeitig mit $\nu + \frac{1}{2}(n+1)$ ein volles Restsystem nach $n$. Daraus ergibt sich nach (1) und (2)

\weq{3}{\Phi\left(u + \frac{1}{n}\right) = -e^{-\frac{4\nu'\pi i}{n}} \Phi(u)}

und ähnlich

\weq{4}{\Phi(u + \omega) = -e^{-n\pi i(2u+\omega)} \Phi(u).}

Daraus ergibt sich, daß die Funktion $\Phi(u)$ eine $t$-Funktion % [sic] The source prints t; theta is intended.
ist mit den Perioden $1/n$, $\omega$, die durch die Bedingungen (3), (4)
bis auf einen von $u$ unabhängigen Faktor bestimmt und durch
eine $\vartheta$-Funktion ausdrückbar ist (§ 19). Man sieht leicht, daß
die Bedingungen (3), (4) durch die Funktion

\[
e^{-4\pi i\nu' u} \vartheta_{11}(nu - 4\nu'\omega, n\omega)
\]
% [sic] The source prints 4 here but 2 in the corresponding argument in (5).

befriedigt sind, so daß sich ergibt:

\weq{5}{\sum^{\nu} e^{\frac{8\nu\nu'\pi i}{n}} \frac{\vartheta_{11}\left(u + \frac{2\nu}{n}\right)}{\vartheta_{g_1,g_2}\left(u + \frac{2\nu}{n}\right)} = Ce^{-4\pi i\nu' u} \frac{\vartheta_{11}(nu - 2\nu'\omega, n\omega)}{\vartheta_{g_1,g_2}(nu, n\omega)}.}

Die Bestimmung der Konstanten $C$ kann man leicht ausführen, wenn man rechts und links mit

\[
\vartheta_{g_1,g_2}(u)
\]

multipliziert und dann

\[
u = \frac{(1-g_2) + (1-g_1)\omega}{2}
\]
%CONT>>
